\chapter{Derivación HMT del parámetro de Barbero--Immirzi y del área mínima: incidencia hexada--octada, canales \(90/120\) y espectro del operador de área}
\label{chap:p3-final:51:barbero_area}
\section{Morfismo de incidencia hexada--octada y conservación de los canales orientados \(90/120\)}

Fijados la dodecada y el alineamiento del
\cref{sec:w12-w24-elevacion}, sea \(H\) una hexada en las coordenadas HMT,
sea \(\ell(H)\subset D\) su imagen y sea
\(O_H=\iota_D(\ell(H))\) su elevación única. Al marcar un punto
y un par de puntos de la hexada residual se obtienen
\[
F_6(H)=\ell(H)\times\binom{\ell(H)}2,
\qquad
|F_6|=6\binom62=90.
\]
Si el punto marcado puede recorrer la octada:
\[
F_8(H)=O_H\times\binom{\ell(H)}2,
\qquad
|F_8|=8\binom62=120.
\]
La inclusión \(\ell(H)\hookrightarrow O_H\) induce el morfismo de incidencia
\[
 \mathcal I_{6\to8}:F_6(H)\hookrightarrow F_8(H).
\]
Al normalizar la diferencia por el factor
elegido \(24\binom62\), se obtiene
\[
\boxed{
\frac{90-120}{24\binom62}
=\frac{6-8}{24}
=-\frac1{12}.}
\]
La igualdad es un defecto normalizado de la interfaz combinatoria
\(6\to8\). El factor \(24\binom62\) forma parte de la definición de esta
normalización; la inclusión por sí sola determina la diferencia \(-30\).
El operador \(\mathcal I_{6\to8}\) tiene como dominio configuraciones
marcadas de incidencia; no es el extractor transversal del estado
dodecafásico ni la familia de momentos angulares.

% Construcción dodecafásica del funcional de retorno.
\section{Cadena fundamental dodecafásica y funcional hexádico--octádico}

La incidencia activa se construye antes de toda realización gravitatoria.
Sea
\[
 H_{\mathrm{act}}=\{1,2,4,5,7,8\},
 \qquad
 \mathcal P_2(H_{\mathrm{act}})
 =\{A\subset H_{\mathrm{act}}:|A|=2\}.
\]
Entonces
\[
 |H_{\mathrm{act}}|=6,
 \qquad
 |\mathcal P_2(H_{\mathrm{act}})|=15,
\]
y las incidencias anterior y posterior a la ampliación poseen cardinales
\[
 \mathcal F_{90}=H_{\mathrm{act}}\times
 \mathcal P_2(H_{\mathrm{act}}),
 \qquad |\mathcal F_{90}|=90,
\]
\[
 \mathcal F_{120}=\{1,\ldots,8\}\times
 \mathcal P_2(H_{\mathrm{act}}),
 \qquad |\mathcal F_{120}|=120.
\]
Estos dos cardinales son salidas de la regla de fase; no son ángulos ni
parámetros recibidos desde la realización física.

Denotemos por \(C_{12}=\mathbb Z/12\mathbb Z\) el registro temporal del TPK y
por \(\partial_{\rm ret}\) su única costura de retorno por vuelta. En el
módulo libre de sucesos orientados se considera la cadena fundamental
\begin{equation}
 c_{12}^{+}
 =\sum_{j\in C_{12}}[j,\mathcal F_{90}]
  -[\partial_{\rm ret},\mathcal F_{120}].
 \label{intsuccpass:eq:c12-fundamental-chain}
\end{equation}
La suma contiene exactamente los doce sectores de la rueda; el segundo
término es único porque una vuelta posee una sola costura de cierre. Su signo
es el signo de frontera. La involución bidireccional envía
\(c_{12}^{+}\) a \(c_{12}^{-}=-c_{12}^{+}\): cambia la orientación de la
cadena, pero no altera sus multiplicidades absolutas ni sus tipos de
incidencia.

Para \(m\in\{90,120\}\), sea
\[
 S_m(q)=\frac{q^m}{1-q^{3m}}.
\]
El lector de retorno \(\mathscr R\) se define sobre los generadores de
\eqref{intsuccpass:eq:c12-fundamental-chain} por
\[
 \mathscr R([j,\mathcal F_{90}])=S_{90},
 \qquad
 \mathscr R([\partial_{\rm ret},\mathcal F_{120}])=S_{120}.
\]
La primera igualdad es constante sobre \(C_{12}\) por covariancia bajo la
rotación de fase; la segunda tiene un dominio distinto y sólo actúa en la
costura ampliada. Por linealidad,
\begin{equation}
 \boxed{\mathscr R(c_{12}^{+})=12S_{90}-S_{120}.}
 \label{intsuccpass:eq:dodecaphase-forces-12minus1}
\end{equation}
Así, los coeficientes \(12:-1\) son la imagen de la cadena fundamental: doce
celdas de fase y una frontera orientada. No se seleccionan mediante el valor
convencional del parámetro de Barbero--Immirzi, ni mediante una comparación
decimal, ni imponiendo previamente el coeficiente de polo que se calculará a
continuación.

\begin{theorem}[Coeficiente de polo de la cadena dodecafásica]
El funcional \(F(q)=\mathscr R(c_{12}^{+})(q)
=12S_{90}(q)-S_{120}(q)\), determinado por los doce sectores
y su costura orientada, tiene coeficiente \(1/24\) en
\((1-q)^{-1}\). Su residuo complejo en \(q=1\) es \(-1/24\).
El coeficiente de polo es una salida del funcional ya construido.
\end{theorem}

\begin{proof}
Para un funcional con polo simple en \(q=1\), escribimos
\[
 p_1(F):=\lim_{q\uparrow1}(1-q)F(q).
\]
La identidad
\[
 1-q^{3m}=(1-q)\sum_{j=0}^{3m-1}q^j
\]
implica
\[
 p_1(S_m)
 =\lim_{q\uparrow1}\frac{q^m}{\sum_{j=0}^{3m-1}q^j}
 =\frac1{3m}.
\]
Por linealidad,
\[
 p_1\bigl(\mathscr R(c_{12}^{+})\bigr)
 =\frac{12}{270}-\frac1{360}
 =\frac{48-3}{1080}=\frac1{24}.
\]
Estas igualdades se efectúan en \(\mathbb Q\). El residuo complejo
es el coeficiente de \((q-1)^{-1}\), de modo que
\[
 \operatorname*{Res}_{q=1}\mathscr R(c_{12}^{+})(q)
 =-p_1\bigl(\mathscr R(c_{12}^{+})\bigr)=-\frac1{24}.
\]
\end{proof}

La ecuación diofántica
\[
 \frac a{270}-\frac b{360}=\frac1{24}
 \quad\Longleftrightarrow\quad
 4a-3b=45
\]
permanece como certificado aritmético posterior. Por sí sola admite la
familia \((a,b)=(12+3t,1+4t)\); la cadena
\eqref{intsuccpass:eq:c12-fundamental-chain} elimina esa ambigüedad antes del cálculo,
pues cambiar \((12,1)\) altera la multiplicidad de la órbita o de su costura
y, por tanto, ya no representa una vuelta fundamental del TPK.

Con los canales conjugados \(q_\pm\) ya generados por la rama angular, se
obtiene finalmente
\[
 \mathscr R(c_{12}^{+})(q_-)-
 \mathscr R(c_{12}^{+})(q_+)
 =12\Delta S_{90}-\Delta S_{120}.
\]
Esta es la entrada interna del funcional hexádico--octádico. Su evaluación
como parámetro de Barbero--Immirzi y su inserción en el operador de área son
publicaciones posteriores: no participan en la selección de
\(c_{12}^{+}\), de \(90/120\), de \(12:-1\) ni de \(1/24\).

\paragraph{Falsadores.}
La construcción queda refutada si el censo de fase no da doce sectores, si
una vuelta contiene más de una costura, si las incidencias no tienen
cardinales \(90\) y \(120\), si la involución no invierte la cadena completa,
o si el coeficiente exacto de \((1-q)^{-1}\) de su imagen difiere de \(1/24\).

\label{mdg:chap:barbero}

La transición de una hexada a una octada combina tres estructuras de tipos
distintos: una incidencia finita, dos series de Lambert asociadas a los
canales \(q_\pm\) y una identificación con el parámetro real de la acción de
Holst. Las tres se articulan mediante la incidencia, el funcional angular y
el diccionario areal hacia la geometría de la conexión de
Ashtekar--Barbero desarrollado en esta misma residencia.

\section{Cardinalidades de incidencia 90 y 120}

Sea \(H\) un conjunto de seis puntos y sea \(O\) un conjunto de ocho puntos
que contiene una copia distinguida de \(H\). Consideremos
\[
 \mathcal F_6=H\times\binom H2,
 \qquad
 \mathcal F_8=O\times\binom H2.
\]
La inclusión distinguida \(H\hookrightarrow O\) induce
\[
 \iota_{6,8}:\mathcal F_6\hookrightarrow\mathcal F_8,
 \qquad
 (h,\{h_1,h_2\})\longmapsto(h,\{h_1,h_2\}).
\]
Este morfismo de incidencias es el antecedente combinatorio de los
cardinales \(90/120\); no es el extractor dodecafásico ni una serie angular.

\begin{theorem}[Cardinalidades de la inclusión]
\[
 |\mathcal F_6|=6\binom62=90,
 \qquad
 |\mathcal F_8|=8\binom62=120.
\]
La diferencia normalizada por veinticuatro coordenadas es
\[
 \frac{90-120}{24\binom62}=-\frac1{12}.
\]
\end{theorem}
\begin{proof}
Cada elemento de la primera coordenada puede combinarse con cualquiera de los
quince pares de \(H\). El principio multiplicativo da las dos cardinalidades;
la última identidad resulta de sustituir \(\binom62=15\).
\end{proof}

\begin{figure}[htbp]
\centering
\begin{tikzpicture}[
  font=\small,
  setnode/.style={circle,draw,minimum size=6mm,inner sep=0pt,
    font=\scriptsize\bfseries},
  box/.style={draw,rounded corners=1.5mm,align=center,
    text width=5.0cm,inner xsep=6pt,inner ysep=5pt},
  arr/.style={-{Stealth[length=2.2mm]},line width=.7pt}]

  % Inclusión geométrica. Las flechas parten del contorno y no atraviesan nodos.
  \node[setnode,draw=hmtblue,fill=hmtlightblue] (h1) at (-2.3,1.25) {1};
  \node[setnode,draw=hmtblue,fill=hmtlightblue] (h2) at (-1.1,1.9) {2};
  \node[setnode,draw=hmtblue,fill=hmtlightblue] (h3) at (.1,1.25) {3};
  \node[setnode,draw=hmtblue,fill=hmtlightblue] (h4) at (.1,-.15) {4};
  \node[setnode,draw=hmtblue,fill=hmtlightblue] (h5) at (-1.1,-.8) {5};
  \node[setnode,draw=hmtblue,fill=hmtlightblue] (h6) at (-2.3,-.15) {6};
  \node[setnode,draw=hmtgreen,fill=hmtlightgreen] (o7) at (-3.2,.55) {7};
  \node[setnode,draw=hmtgreen,fill=hmtlightgreen] (o8) at (1.0,.55) {8};
  \node[draw=hmtblue,dashed,rounded corners=4mm,
    fit=(h1)(h2)(h3)(h4)(h5)(h6),inner sep=4mm,
    label={[hmtblue,font=\bfseries]above:{hexada \(H\), \(|H|=6\)}}] (H) {};
  \node[draw=hmtgreen,line width=.8pt,rounded corners=6mm,
    fit=(H)(o7)(o8),inner sep=5mm,
    label={[hmtgreen,font=\bfseries]below:{octada \(O\), \(|O|=8\), \(H\subset O\)}}] (O) {};

  \node[box,draw=hmtblue,fill=hmtlightblue] (f6) at (5.3,1.35)
    {Incidencias definidas sobre \(H\)\\[1mm]
     \(\mathcal F_6=H\times\binom H2\),\quad
     \(|\mathcal F_6|=6\binom62=90\)};
  \node[box,draw=hmtgreen,fill=hmtlightgreen] (f8) at (5.3,-1.0)
    {Extensión de incidencias a \(O\)\\[1mm]
     \(\mathcal F_8=O\times\binom H2\),\quad
     \(|\mathcal F_8|=8\binom62=120\)};
  \draw[arr,hmtblue] (O.east) -- ++(.45,0) |- (f6.west);
  \draw[arr,hmtgreen] (O.east) -- ++(.45,0) |- (f8.west);
  \draw[arr,hmtgold] (f6.south) -- node[right,align=left,font=\footnotesize]
    {inclusión \(\iota_{6,8}\)} (f8.north);

  \node[box,draw=hmtgold,fill=hmtlightgold,text width=10.9cm]
    (def) at (1.45,-3.80)
    {Defecto combinatorio orientado\\[-.2mm]
     con la normalización declarada:\\[1mm]
     \(\displaystyle
       \frac{|\mathcal F_6|-|\mathcal F_8|}{24\binom62}
       =\frac{90-120}{24\cdot15}=-\frac1{12}.\)};
  \draw[arr,hmtgold] (f8.south) -- (f8.south |- def.north);

  \node[align=center,text=hmtgray,font=\footnotesize,text width=11.2cm]
    at (1.45,-5.50)
    {Esta identidad pertenece al mapa de incidencias. No identifica el
     extractor dodecafásico, el semigrupo de torres ni el funcional angular.};
\end{tikzpicture}

\caption{Inclusión de un subconjunto distinguido de seis elementos en uno de
ocho elementos y conteos de incidencia. Los cardinales \(90\) y \(120\) no fusionan el extractor
dodecafásico, el funcional angular ni el semigrupo de torres.}
\label{mdg:fig:incidencia-seis-ocho-rev9}
\end{figure}

Los índices iniciales de la jerarquía angular quedan así asociados a dos
espacios de incidencia: 90 configuraciones sobre la hexada y 120 sobre la
octada ambiente. El cociente \(-1/12\) es una identidad combinatoria de esta
normalización.

\section{Series de Lambert asociadas a los 2 sectores de incidencia}

Para \(0<q<1\), definimos simultáneamente
\[
 S_{90}(q)=\frac{q^{90}}{1-q^{270}},
 \qquad
 S_{120}(q)=\frac{q^{120}}{1-q^{360}},
 \qquad
 \Delta S_m=S_m(q_-)-S_m(q_+).
\]
Las expansiones geométricas
\[
 \Delta S_{90}=\sum_{j\geq0}s_{90+270j},
 \qquad
 \Delta S_{120}=\sum_{j\geq0}s_{120+360j}
\]
sitúan ambos términos dentro del semigrupo global del
\cref{mdg:chap:torres}.

El funcional angular asociado a la transición \(6\to8\) es
\[
 K_{6\to8}=12\Delta S_{90}-\Delta S_{120},
\]
y su precedencia ha quedado fijada por la cadena fundamental dodecafásica:
los doce sectores y la única costura orientada producen primero el peso
\(12:-1\). El coeficiente de \((1-q)^{-1}\), igual a \(1/24\), se calcula después como salida exacta;
la ecuación diofántica que sigue certifica el par ya generado y no lo
selecciona.
La coordenada angular resultante es
\[
\Gamma_{6\to8}
=\frac{\pi A\Cstar}{180}+K_{6\to8}.
\]
En el punto HMT,
\[
\begin{aligned}
 \Delta S_{90}&=2.95169439297457621139847987861\times10^{-4},\\
 \Delta S_{120}&=1.96835112502951720093738706217\times10^{-5},\\
 \frac{\pi A\Cstar}{180}&=0.270485322934140078465586458790\ldots,
\end{aligned}
\]
y, por tanto,
\[
 \boxed{\Gamma_{6\to8}
 =0.27400767269445927474725526077398041940\ldots.}
\]
Las tablas calculadas con la coordenada angular conjugada \(C^*\) redondeada a quince decimales dan
\(0.2740076726944593\) en aritmética de doble precisión; el valor anterior
es la evaluación multiprecisión de la fase regional completa.

\section{Certificado diofántico posterior de los pesos}

La identidad
\[
 \lim_{q\to1^-}(1-q)\frac{q^m}{1-q^n}=\frac1n
\]
asocia a la combinación \(aS_{90}-bS_{120}\), antes de su evaluación
en los canales conjugados, el coeficiente de \((1-q)^{-1}\)
\[
 \frac a{270}-\frac b{360}.
\]

\begin{theorem}[Pesos enteros con coeficiente de polo \texorpdfstring{\(1/24\)}{1/24}]
Las soluciones enteras positivas de
\[
 \frac a{270}-\frac b{360}=\frac1{24}
\]
son
\[
 (a,b)=(12+3t,1+4t),\qquad t\in\ZZ_{\geq0}.
\]
El par \((12,1)\) es la única solución mínima coordenada a coordenada y la
única de norma \(\ell^1\) mínima.
\end{theorem}
\begin{proof}
Multiplicar por \(1080\) produce \(4a-3b=45\). Módulo tres se obtiene
\(a\equiv0\pmod3\); escribiendo \(a=3u\), resulta \(4u-b=15\). La
positividad implica \(u=4+t\), \(b=1+4t\), con \(t\geq0\). Por tanto,
\(a+b=13+7t\), lo que demuestra la minimalidad.
\end{proof}

La primitividad \(\gcd(a,b)=1\) no selecciona por sí sola el par mínimo, pues
existen otras soluciones primitivas. La vuelta fundamental del TPK ya ha
producido el peso \(12:-1\); el coeficiente de polo \(1/24\), la positividad y la
minimalidad forman su certificado aritmético posterior. Ninguno de estos
objetos se deduce de la evaluación decimal de \(\Gamma_{6\to8}\).

\newpage
\section{Monotonía del funcional respecto de la coordenada angular conjugada}

Para \(m\in\{90,120\}\),
\[
 S_m'(q)=\frac{m q^{m-1}(1+2q^{3m})}{(1-q^{3m})^2}>0.
\]
Además,
\[
 \frac{\dd q_-}{\dd C}=\frac\pi{180}q_-,
 \qquad
 \frac{\dd q_+}{\dd C}=-\frac\pi{180}q_+.
\]

\begin{proposition}[Monotonía estricta]
Para \(1^\circ\leq C\leq3^\circ\), la función
\(C\mapsto\Gamma_{6\to8}(A,C)\) es estrictamente creciente. En consecuencia,
todo valor de su imagen tiene a lo sumo una preimagen en ese intervalo.
\end{proposition}
\begin{proof}
La derivada de cada diferencia \(\Delta S_m\) es positiva. En el intervalo
indicado se obtiene
\(\Delta S'_{120}<5.35\times10^{-4}\), mientras que
\(\pi A/180>0.1273\). Por tanto,
\(\Gamma'(C)>0.1267\). En \(C=\Cstar\),
\[
 \Gamma'(\Cstar)=0.1328995105618907\ldots.
\]
\end{proof}

La inyectividad local excluye una segunda coordenada angular próxima con la
misma evaluación del funcional. Este control no construye ni selecciona
\(C^*\): la coordenada angular conjugada ya ha sido obtenida en la rama regional anterior.
La proposición sólo cuantifica la sensibilidad de la evaluación descendente
una vez fijados \(A\) y \(C^*\); no autoriza a invertir
\(\Gamma_{6\to8}\) ni un valor físico de \(\gamma_{\mathrm{BI}}\) para
redefinirlos.

\label{ch:modular}

\section{Inclusi\'on hexada--octada y realizaciones funcionales asociadas}

Fijemos un conjunto \(H\) de seis puntos y un conjunto \(O\) de ocho
puntos con una inclusi\'on distinguida \(j:H\hookrightarrow O\). En la
residencia de Witt, \(H\) es una hexada de un dodecad y \(O\) su octada
elevada; en este cap\'itulo s\'olo se usa la inclusi\'on ya fijada, no el
valor terminal de Barbero--Immirzi. Escribimos
\begin{equation}
\label{eq:modular-pairs-H}
\binom H2:=\{P\subset H:|P|=2\},
\qquad \left|\binom H2\right|=\binom62=15,
\end{equation}
y definimos dos espacios de incidencias con tipos distintos:
\begin{equation}
\label{eq:modular-F6-F8}
\mathcal F_6:=H\times\binom H2,
\qquad
\mathcal F_8:=O\times\binom H2.
\end{equation}
La inclusi\'on de incidencias es
\begin{equation}
\label{eq:modular-iota68}
\iota_{6,8}:\mathcal F_6\hookrightarrow\mathcal F_8,
\qquad
(h,P)\longmapsto(j(h),P).
\end{equation}
Es inyectiva porque \(j\) lo es y conserva el par marcado \(P\).

\begin{theorem}[Conteos de incidencia y defecto orientado]
\label{thm:modular-incidence-90-120}
Los dos espacios de \cref{eq:modular-F6-F8} satisfacen
\begin{equation}
\label{eq:modular-counts-90-120}
|\mathcal F_6|=6\binom62=90,
\qquad
|\mathcal F_8|=8\binom62=120,
\qquad
|\mathcal F_8\setminus\iota_{6,8}(\mathcal F_6)|=30.
\end{equation}
La normalizaci\'on por las veinticuatro coordenadas del par
dodecad--veinticuatro y por los quince pares de \(H\) produce
\begin{equation}
\label{eq:modular-incidence-defect}
\boxed{
\delta_{6,8}^{\rm inc}
=\frac{|\mathcal F_6|-|\mathcal F_8|}
{24\binom62}=-\frac1{12}.}
\end{equation}
La misma fracci\'on se obtiene mediante el funcional rec\'iproco asociado
a la memoria incorporada,
\begin{equation}
\label{eq:modular-reciprocal-defect}
\boxed{
\delta_{6,8}^{\rm rec}
=\frac{30}{120}-\frac{30}{90}=-\frac1{12}.}
\end{equation}
Son dos realizaciones exactas del mismo defecto orientado, pero tienen
dominios distintos y no se identifican con \(\zeta(-1)\) sin un
entrelazador adicional.
\end{theorem}

\begin{proof}
El principio multiplicativo aplicado a
\cref{eq:modular-F6-F8} da \(6\cdot15=90\) y \(8\cdot15=120\).
La imagen de \(\iota_{6,8}\) tiene noventa elementos, de modo que el
complemento tiene treinta. Sustituir en
\cref{eq:modular-incidence-defect} da
\(-30/(24\cdot15)=-1/12\), mientras
\cref{eq:modular-reciprocal-defect} es
\(1/4-1/3=-1/12\).
\end{proof}

La transici\'on hexada--octada admite dos realizaciones funcionales que no
se deducen mutuamente. El diagrama algebraico, formulado sobre los espacios
que justifican los conteos, es
\begin{equation}
\label{eq:modular-68-two-readers}
\begin{aligned}
(H\xhookrightarrow{\,j\,}O)
&\longmapsto
\left(|\mathcal F_6|,|\mathcal F_8|\right)=(90,120)
\longmapsto\Gamma_{6\to8}=\gammaH,\\
(h_6,o_8,t)=(6,8,3)
&\longmapsto-\frac{o_8-h_6}{t}=-\frac23,
\end{aligned}
\end{equation}
donde \(-2/3\) es el coeficiente de la tercera coordenada de la fila
\(13\) del constructor CKM. La primera realizaci\'on transforma la
incidencia en \`area y memoria modular; la segunda transforma la misma
transici\'on en par\'ametros de mezcla angular.
No se postula \(\boldsymbol\theta_{\rm CKM}=f(\gammaH)\): ambas salidas
conservan un antecedente com\'un.

El morfismo que determina los sectores \(90/120\) es
\(\iota_{6,8}\). La magnitud \(\gammaH\) ya ha sido determinada internamente
por la construcci\'on hexada--octada y se reutiliza aqu\'i sin reconstruirla a
partir de su representaci\'on decimal. Las
consecuencias operatorias comienzan s\'olo despu\'es de fijar la escala
elemental \(a_0\), que se define a continuaci\'on sin usar entrop\'ia ni el
producto terminal \(\gammaH a_0\).

\section{Correspondencia entre el interior non\'adico y los grados de libertad de borde}

Sea \(\mathbb T_3=\{0,1,2\}\). Todo \(x\in\mathbb Z/9\mathbb Z\) posee
una \`unica escritura de d\'igitos
\(x=x_0+3x_1\), con \(x_i\in\mathbb T_3\), y todo
\(y\in\mathbb Z/27\mathbb Z\) una \`unica escritura
\(y=y_0+3y_1+9y_2\). Por ello
\begin{equation}
\label{eq:modular-bulk-boundary-729}
\mathcal B=(\mathbb Z/9\mathbb Z)^3,
\qquad
\mathcal B_\partial=(\mathbb Z/27\mathbb Z)^2,
\qquad
|\mathcal B|=|\mathcal B_\partial|=3^6=729.
\end{equation}
La igualdad de cardinales se acompa\~na del mapa expl\'icito. Si
\(x_j=r_{2j-1}+3r_{2j}\) para \(j=1,2,3\), definimos
\begin{equation}
\label{eq:modular-beta729}
\beta_{729}(x_1,x_2,x_3)
=\bigl(r_1+3r_2+9r_3,\;r_4+3r_5+9r_6\bigr).
\end{equation}

\begin{proposition}[Bijecci\'on geneal\'ogica \(729\leftrightarrow729\)]
\label{prop:modular-beta729}
El mapa \(\beta_{729}:\mathcal B\to\mathcal B_\partial\) es biyectivo y
preserva la palabra ordenada de seis trits. No es, en general, un
isomorfismo de grupos aditivos: reagrupar los d\'igitos cambia d\'onde se
registra el acarreo.
\end{proposition}

\begin{proof}
Extraer los seis d\'igitos de tres coordenadas non\'adicas y agruparlos en
dos ternas es una permutaci\'on de coordenadas de \(\mathbb T_3^6\). La
descomposici\'on posicional es \`unica en ambos lados, por lo que la
operaci\'on tiene inversa. La advertencia aditiva se comprueba, por ejemplo,
al sumar palabras que generan acarreo entre \(r_3\) y \(r_4\): esa
frontera separa las dos coordenadas de orden veintisiete, pero atraviesa
dos coordenadas distintas en el empaquetado non\'adico.
\end{proof}

En el toro de borde \(\mathcal B_\partial\), tomamos el bloque de v\'ertices
\begin{equation}
\label{eq:modular-block-A}
A=\{0,1,\ldots,8\}\times\{0,1,\ldots,8\}.
\end{equation}
Sus enlaces orientados que cruzan el borde se separan por direcci\'on:
\begin{align}
\partial A_{90}
&=\{((-1,j),(0,j)),((8,j),(9,j)):0\leq j\leq8\},
\label{eq:modular-boundary90}\\
\partial A_{120}
&=\{((i,-1),(i,0)),((i,8),(i,9)):0\leq i\leq8\},
\label{eq:modular-boundary120}
\end{align}
donde todas las coordenadas se leen m\'odulo veintisiete. As\'i
\begin{equation}
\label{eq:modular-boundary-counts}
\partial A=\partial A_{90}\sqcup\partial A_{120},
\qquad
|\partial A_{90}|=|\partial A_{120}|=18.
\end{equation}
Los sub\'indices \(90,120\) etiquetan la procedencia de canal; no dicen que
el borde tenga noventa o ciento veinte enlaces.

\begin{theorem}[Ley areal del corte tri\'adico]
\label{thm:modular-triadic-cut}
En el grafo c\'ubico \(N\times N\times N\) con capacidad unitaria, todo
corte entre dos caras opuestas posee capacidad al menos \(N^2\), y una
capa transversal alcanza esa cota. Para \(N=9\),
\begin{equation}
\label{eq:modular-mincut-entropy}
\mathcal A_{\min}=81,
\qquad
S_{\rm tri}=81\log3.
\end{equation}
\end{theorem}

\begin{proof}
Hay \(N^2\) rectas de aristas paralelas a la direcci\'on normal, disjuntas
dos a dos. Todo corte debe interceptar cada una, de modo que su capacidad
es al menos \(N^2\). Una sola capa transversal contiene exactamente
\(N^2\) aristas y realiza el m\'inimo. Si cada enlace cortado porta un trit
m\'aximamente mezclado, su entrop\'ia es \(\log3\); la aditividad sobre los
ochenta y un enlaces da \cref{eq:modular-mincut-entropy}.
\end{proof}

La entrop\'ia \(S_{\rm tri}\) mide capacidad del corte c\'ubico. La
entrop\'ia modular de los treinta y seis enlaces de
\cref{eq:modular-boundary-counts} mide otro estado y no debe igualarse con
ella.

\section{Entrelazador colectivo de incidencia entre interior y borde}

En la suma
\(\ell^2(\mathcal F_6)\oplus\ell^2(\mathcal F_8)\), definimos
\begin{equation}
\label{eq:modular-u90-u120}
u_{90}=\frac1{\sqrt{90}}\sum_{f\in\mathcal F_6}(\delta_f,0),
\qquad
u_{120}=\frac1{\sqrt{120}}\sum_{f\in\mathcal F_8}(0,\delta_f).
\end{equation}
En
\(\ell^2(\partial A_{90})\oplus\ell^2(\partial A_{120})\), sean
\begin{equation}
\label{eq:modular-v90-v120}
v_{90}=\frac1{\sqrt{18}}\sum_{e\in\partial A_{90}}(\delta_e,0),
\qquad
v_{120}=\frac1{\sqrt{18}}\sum_{e\in\partial A_{120}}(0,\delta_e).
\end{equation}
Las parejas \((u_{90},u_{120})\) y \((v_{90},v_{120})\) son ortonormales.
Por tanto existe una \`unica isometr\'ia entre los subespacios colectivos
que satisface
\begin{equation}
\label{eq:modular-J68}
\mathcal J_{6\to8}u_{90}=v_{90},
\qquad
\mathcal J_{6\to8}u_{120}=v_{120}.
\end{equation}
Definamos
\begin{align}
D_{\rm inc}
&=90|u_{90}\rangle\langle u_{90}|
+120|u_{120}\rangle\langle u_{120}|,
\label{eq:modular-Dinc}\\
D_{\rm area}
&=90|v_{90}\rangle\langle v_{90}|
+120|v_{120}\rangle\langle v_{120}|.
\label{eq:modular-Darea}
\end{align}

\begin{theorem}[Transporte exacto de las dos etiquetas]
\label{thm:modular-J68}
La aplicaci\'on \(\mathcal J_{6\to8}\) es unitaria entre los dos planos
colectivos y entrelaza los operadores de incidencia y \`area:
\begin{equation}
\label{eq:modular-J-intertwining}
\boxed{
\mathcal J_{6\to8}D_{\rm inc}
=D_{\rm area}\mathcal J_{6\to8}.}
\end{equation}
Por tanto, toda combinaci\'on polin\'omica de las etiquetas \(90,120\), en
particular la combinaci\'on primitiva \(12:-1\), se transporta con los
mismos coeficientes.
\end{theorem}

\begin{proof}
La ortonormalidad muestra que \(\mathcal J_{6\to8}\) env\'ia una base
ortonormal en otra. En \(u_{90}\) ambos lados de
\cref{eq:modular-J-intertwining} valen \(90v_{90}\); en \(u_{120}\),
ambos valen \(120v_{120}\). La igualdad en una base prueba la identidad.
El c\'alculo funcional polin\'omico conserva el entrelazamiento.
\end{proof}

Este mapa no pretende inyectar noventa incidencias individuales en
dieciocho enlaces: transporta exactamente los dos modos uniformes y sus
etiquetas espectrales. Extenderlo a fluctuaciones no uniformes exigir\'ia
un nuevo mapa de incidencia; el operador modular definido aqu\'i s\'olo emplea
el subespacio colectivo que acaba de construirse.

\section{Determinación de la escala elemental del espectro del operador de área}

Sea \(C^*\) la coordenada angular conjugada, expresada en grados, y fijemos
\begin{equation}
\label{eq:modular-theta-clock-a0}
\theta_{\rm clk}=\frac{C^*}{6}\frac{\pi}{180},
\qquad
\boxed{
a_0=\frac{1+2\cos(10^\circ)}{3}\,\theta_{\rm clk}.}
\end{equation}
El factor que precede a \(\theta_{\rm clk}\) no es una correcci\'on
decimal. Es el car\'acter real normalizado del triplete de fases
\begin{equation}
\label{eq:modular-c10-character}
\frac13\Tr\operatorname{diag}
\left(1,\ee^{\ii\pi/18},\ee^{-\ii\pi/18}\right)
=\frac{1+2\cos(10^\circ)}3.
\end{equation}
As\'i, \(a_0\) es la media espectral de la estructura angular sobre el trit
orientado. Es positiva, se fija antes del operador de \`area y no se
selecciona a partir de la entrop\'ia. La identificaci\'on posterior
\(\gammaH=\Gamma_{6\to8}\) realiza la salida HMT en la conexi\'on de
Ashtekar--Barbero \cite{Barbero,Immirzi}; no reabre su derivaci\'on.

\begin{remark}[Alcance del s\'imbolo \(a_0\)]
En este cap\'itulo, \(a_0\equiv a_0^{\rm area}\) denota exclusivamente la
normalizaci\'on de \cref{eq:modular-theta-clock-a0}. No es el radio de Bohr
ni el factor de escala cosmol\'ogico inicial, aunque esos objetos se escriban
tradicionalmente con el mismo s\'imbolo. Ninguna identidad de este cap\'itulo
permite sustituir uno por otro.
\end{remark}

\section{Ley areal del corte triádico: capacidad \(81\log 3\), operador de área, espectro y entrelazador colectivo de los canales \(90/120\)}
\label{hap:sec:area-entrelazamiento-barbero-rev11}
\index{Barbero--Immirzi!operador de área}
\index{entropía modular!Hamiltoniano modular}

La identificación HMT del funcional \(\Gamma_{6\to8}\) con la coordenada
de Barbero--Immirzi se realiza sobre el mismo borde triádico que soporta la
incidencia (90/120). El interior y su hoja de borde poseen igual cardinalidad,
\[
 \mathcal B=(\mathbb Z/9\mathbb Z)^3,
 \quad |\mathcal B|=729,
 \qquad
 \partial\mathcal B=(\mathbb Z/27\mathbb Z)^2,
 \quad |\partial\mathcal B|=729,
\]
mediante el reagrupamiento de seis trits en tres coordenadas nonádicas o dos
coordenadas de orden veintisiete.

\begin{theorem}[Ley areal del corte triádico]
\label{hap:thm:ley-areal-corte-triadico-rev11}
En el grafo cúbico \(N\times N\times N\) con capacidad unitaria, el corte
mínimo entre dos caras opuestas tiene capacidad \(N^2\). Para \(N=9\),
\[
 \mathcal A_{\min}=81,
 \qquad S_{\rm tri}=\mathcal A_{\min}\log3=81\log3.
\]
\end{theorem}
\begin{proof}
Las \(N^2\) rectas paralelas a la dirección normal son caminos disjuntos en
aristas, por lo que todo corte las intercepta. Una capa transversal contiene
exactamente \(N^2\) aristas y alcanza la cota. Cada enlace triádico aporta
\(\log3\) a la entropía de un estado máximamente mezclado.
\end{proof}

En el toro de borde \(27\times27\), sea \(A\) el bloque canónico
\(9\times9\). Su frontera se descompone de manera exacta en
\begin{equation}
 \partial A=\partial A_{90}\sqcup\partial A_{120},
 \qquad |\partial A_{90}|=|\partial A_{120}|=18.
 \label{hap:eq:corte-90-120-rev11}
\end{equation}
Los enlaces horizontales representan el canal (90), y los verticales, el
canal (120). Esta descomposición convierte la pareja de incidencias en una
descomposición operatoria del borde.

La entropía areal \(S_{\rm tri}=81\log3\) pertenece al estado máximamente
mezclado de los enlaces del corte. La entropía bicapa
\(S_{\rm bi}(y)\) de la \cref{hap:def:entropia-bicapa-rev6} pertenece a la
distribución normalizada entre las dos hojas; ambos objetos conservan
dominios distintos.

Sea \(N=\operatorname{diag}(0,1,2)\) el operador número de un enlace trítico.
La realización areal declara la escala positiva
\[
 \theta_{\rm clk}=\frac{C^*}{6}\frac{\pi}{180},
 \qquad
 a_0=\frac{1+2\cos(10^\circ)}{3}\,\theta_{\rm clk},
 \qquad \gamma_{\rm HMT}=\Gamma_{6\to8}.
\]
La fórmula de \(a_0\) constituye la estructura de partida de esta
realización; fijada esa escala, el operador y su espectro se deducen
exactamente.

\begin{definition}[Operador de área HMT]
Sobre la hoja de borde
\(\mathcal H_{\partial}=\bigotimes_{e\in\partial A}\mathbb C^3\), se define
\begin{equation}
 \widehat{\mathcal A}_{\rm HMT}
 =\gamma_{\rm HMT}a_0\sum_{e\in\partial A}N_e.
 \label{hap:eq:operador-area-hmt-rev11}
\end{equation}
Su espectro en el corte canónico es
\[
 \operatorname{Spec}\widehat{\mathcal A}_{\rm HMT}
 =\{n\gamma_{\rm HMT}a_0:0\leq n\leq72\},
\]
con multiplicidades dadas por los coeficientes de
\((1+z+z^2)^{36}\). La prolongación por cortes compatibles produce la
familia protegida \(\mathcal A_n^{\rm prot}=n\gamma_{\rm HMT}a_0\).
\end{definition}

El entrelazador que fija la normalización se construye en el subespacio de
canales colectivos. Sean \(u_{90},u_{120}\) los vectores unitarios uniformes
de las incidencias hexádica y octádica, y sean \(v_{90},v_{120}\) los
vectores unitarios uniformes de los dos conjuntos de enlaces de
\eqref{hap:eq:corte-90-120-rev11}. El mapa
\begin{equation}
 \begin{aligned}
 \mathcal J_{6\to8}:{}
 &\operatorname{span}\{u_{90},u_{120}\}
 \longrightarrow
 \operatorname{span}\{v_{90},v_{120}\},\\
 &\mathcal J_{6\to8}u_{90}=v_{90},
 \qquad
 \mathcal J_{6\to8}u_{120}=v_{120}
 \end{aligned}
 \label{hap:eq:entrelazador-area-barbero-rev11}
\end{equation}
es una isometría entre estos subespacios colectivos. Si
\[
 D_{\rm inc}=90|u_{90}\rangle\langle u_{90}|
 +120|u_{120}\rangle\langle u_{120}|,
 \qquad
 D_{\rm area}=90|v_{90}\rangle\langle v_{90}|
 +120|v_{120}\rangle\langle v_{120}|,
\]
entonces
\[
 \boxed{\mathcal J_{6\to8}D_{\rm inc}
 =D_{\rm area}\mathcal J_{6\to8}.}
\]
En particular, la combinación primitiva \(12:-1\) actúa
con los mismos coeficientes antes y después de \(\mathcal J_{6\to8}\). Esta
igualdad fija el diccionario de normalización que inserta
\(\gamma_{\rm HMT}=\Gamma_{6\to8}\) en la conexión de
Ashtekar--Barbero.


El paso de hexada a octada es una transformación de la estructura de incidencia.\par

Al pasar de \(6\) a \(8\), cambian los pares disponibles, las relaciones internas y la forma en que la orientación puede conservarse. Esa transición genera un defecto orientado. La variación cardinal realiza un giro en el modo de organización.\par

La misma inversión de hoja se expresa mediante\par

\[
C^*\longmapsto -C^*.
\]

Pero distintos lectores responden de formas distintas a esa inversión.\par

El funcional hexada–octada que conduce al parámetro de Barbero–Immirzi es impar: al invertir la hoja, cambia de signo. La orientación geométrica se invierte.\par

La entropía bicapa, en cambio, es par: atribuye la misma entropía a \(C^*\) y \(-C^*\) cuando se observa la distribución de probabilidades. La información total puede permanecer igual aunque la orientación se haya invertido.\par

El espacio CKM realiza una tercera respuesta. La inversión induce sobre los tres ángulos una reflexión racional exacta:\par

\[
R_{\rm CKM}
=
M_{\rm CKM}
\operatorname{diag}(1,-1,1)
M_{\rm CKM}^{-1},
\]

con\par

\[
R_{\rm CKM}^2=I,
\qquad
\det R_{\rm CKM}=-1.
\]

Esto significa que el cambio de hoja se transporta al espacio de mezcla como una verdadera reflexión. Existe una dirección impar de rango uno y un plano que permanece fijo.\par

El contraángulo \(C^*\) es precisamente la coordenada que cambia. La fase construida desde \(A\) permanece invariante. La transformación separa así dos memorias:\par

\begin{itemize}[leftmargin=*,itemsep=0.35em]
\item una memoria de orientación;
\item una memoria de fase.
\end{itemize}

Esto ofrece una lectura muy profunda de CKM. Los ángulos de mezcla son coordenadas de un espacio sobre el que actúa la misma involución que ya organizaba la transición hexada–octada.\par

El giro \(6\to8\) suministra la operación de orientación que CKM realiza en su propio codominio; la causalidad física pertenece al transductor correspondiente.\par

Tenemos entonces tres publicaciones hermanas del mismo tránsito:\par

\begin{enumerate}[leftmargin=*,itemsep=0.65em]
\item \textbf{Publicación geométrica.}
El parámetro de Barbero–Immirzi orienta el cuanto de área.\par

\item \textbf{Publicación informacional.}
El estado reducido y su Hamiltoniano modular miden qué información de esa orientación permanece en el borde.\par

\item \textbf{Publicación de mezcla.}
La reflexión CKM transforma las coordenadas de mezcla y aísla su componente impar.\par
\end{enumerate}

\(\gamma_{\rm BI}\), la entropía y los ángulos CKM permanecen como realizaciones distintas de una bisagra genealógica común. La unidad reside en el antecedente anterior a la ramificación. Comparten la bisagra, pero cada uno conserva una parte distinta de ella.\par

Ésta es una de las expresiones más bellas de la metodología HMT: la unidad fundamental se recupera remontando las genealogías hasta encontrar la operación común que genera resultados terminales diferentes.\par

La gravedad conserva el giro como orientación del área.\par
La información conserva el giro como memoria del estado.\par
La mezcla conserva el giro como reflexión de coordenadas.\par

La misma transformación «respira» de tres maneras.\par


\section{Separación causal entre el espectro de área, la memoria modular y la realización Cartan--Holst}
La derivación de $\gamma_{\mathrm{BI}}$, del área mínima y del espectro del operador de área concluye en este capítulo. La reconstrucción conjunta de ocupación y procedencia mediante el Hamiltoniano modular de borde, y la realización Palatini--Cartan--Holst, conservan dominios, operadores y pruebas propios en la Parte IV.
